\begingroup
\selectlanguage{english}
\noindent\textbf{Chapter XIX: factorization, diagram categories, and the existence of derived Kan extensions.}

\smallskip
\noindent\emph{Commentary and mathematical audit produced by GPT-6 Astra
(OpenAI, ChatGPT Pro mode). The source is Grothendieck's text in the electronic
edition deciphered and transcribed by M.~K\"unzer and edited by M.~K\"unzer,
J.~Malgoire, and G.~Maltsiniotis. This delivery does not certify any subsequent
assembly or editing with Claude (Anthropic).}

\medskip
Grothendieck has reached the point at which his formalism must acquire its
examples. Given a category $\mathcal M$ and a class $W$ of weak equivalences,
the intended coefficient system is
\[
 \mathbb D_{\mathcal M,W}(I)
   =\operatorname{Fun}(I^{\mathrm{op}},\mathcal M)[W(I)^{-1}],
 \qquad u^*:\mathbb D(J)\longrightarrow\mathbb D(I).
\]
The problem is to construct both adjoints to restriction, prove their pointwise
formulas, and explain which hypotheses on the models are doing the work. The
chapter's fifty manuscript pages therefore move from lifting and factorization
to model structures on diagrams, and then return to the places where that
method appears to stop. In the edition the chapter occupies printed pages 1--26,
after two preliminary leaves.

The historical motivation precedes the terminology of this chapter. In
\emph{Pursuing Stacks}, \S69, Grothendieck relates the passage from a single
derived category to a system of derived diagram categories to the inadequacy
of the former for the variance and trace constructions arising in SGA~5.
His letter to Thomason of 2 April 1991 connects the long investigation of
accessible categories in Chapter XVIII with the foundations needed for this
coefficient formalism \citep{Grothendieck1983,Grothendieck1991Letter}. Thus
accessibility enters here for a concrete reason: the homotopical operations
must be constructed using sets of generators and controlled colimits before
they can be recognized intrinsically after localization.

\noindent\textbf{The two levels of closure.}
The opening pages isolate the relation $i\mathbin{\perp}p$ of admitting a
diagonal in every lifting square. It gives an antitone correspondence between
classes of arrows, hence the closures
${}^{\perp}(S^{\perp})$ and $({}^{\perp}T)^{\perp}$. Pushouts, transfinite
compositions, and retracts describe a second closure process. The small
object argument identifies the two when a set of arrows generates the
lifting condition and their domains have the required smallness. The
retract argument is the decisive step: factor an arrow already having the
left lifting property, then use its lift against the right factor to express
it as a retract of the constructed cellular arrow. This is why the distinction
between an orthogonal pair and a factorizing pair must remain visible.
The source's initial identification of separate closedness of two orthogonal
classes with mutual determination of the classes is too strong; the
factorization hypothesis supplies the missing implication later.

Grothendieck's word ``special'' must also retain its stated meaning. His
ambient category has small colimits and each object is presentable at some
cardinal; his special left classes have a small set of lifting generators.
This does not itself supply the small generating subcategory required for
local presentability. Modern statements about combinatorial model categories
may consequently be used only after the additional ambient hypothesis has
been made explicit.

\noindent\textbf{The asymmetry of projective and injective diagrams.}
For contravariant diagrams, evaluation at $i$ has a left adjoint
$i_!A(j)=\coprod_{\operatorname{Hom}_I(j,i)}A$. Applying these adjoints to
generating cofibrations and generating trivial cofibrations constructs the
projective structure. Its fibrations and weak equivalences are pointwise;
its cofibrations encode compatibility with the arrows of the indexing
category. Theorem 2.3 then derives $u_!$, and Proposition 1.12 supplies the
restriction property needed to compute its values on comma categories.
In the variance used here the formulas are
\[
 (\mathbf L u_!X)_j\simeq
   \operatorname*{hocolim}_{(j\downarrow u)^{\mathrm{op}}}X,
 \qquad
 (\mathbf R u_*X)_j\simeq
   \operatorname*{holim}_{(u\downarrow j)^{\mathrm{op}}}X.
\]
Writing the opposite categories is necessary: the manuscript's diagrams
are presheaves, whereas much modern derivator notation uses covariant
diagrams.

For the injective structure the cofibrations are pointwise, and it is the
right class whose compatibility is no longer pointwise. One cannot obtain
small generators merely by placing the old generators at every component.
Question 1.11 has a positive answer for locally presentable ambient
categories, by the generation theorem of \citet[Lemma A.2.8.3]{Lurie2009HTT},
but its final suggested choice of generators already fails for sets with a
nontrivial group action. The commentary gives the sharper elementary
criterion for when the sole generator $\varnothing\to 1$ can be used unchanged
on a diagram category, with a complete proof and an explicit statement of
the literature checked.

Theorems 2.6 and 2.9 secure the injective construction in Grothendieck's
``very special'' situation, in which the cofibrations are all monomorphisms.
The geometric conditions on unions and pushouts, together with accessibility
of $W$, are the mechanism behind the recognition theorem. Here two-out-of-three
does essential work: after building a bounded trivial subextension, it keeps
the remaining extension in $W$. This is a substantially stronger situation
than an arbitrary accessible class of monomorphisms. The proof recalled
from Chapter XV must be read through its corrected argument in \S7, not
through the explicitly abandoned general claim in \S6. The cardinal-stage
condition printed with a projective-limit arrow in XIX is identified in the
apparatus as a source slip: the cited formulations require closure under
the corresponding small inductive limits.

The latching interpretation in section 3 returns to the same care about
testing classes. Proposition 3.5 needs the canonical right class, which the
source explicitly chooses in (3.5.2); an arbitrary presentation by right
tests can give a false equality. The commentary supplies the counterexample
and the valid comparison. It also proves the classwise cellular description
conjectured in (3.5.3) for noetherian posets, with the required colimits stated,
and supplies the short lifting lemma implicit in Proposition 3.7.

\noindent\textbf{What later existence theorems change.}
The existence of a convenient diagram model structure and the existence of
derived Kan extensions are different assertions. Cisinski proves the latter
for an arbitrary Quillen model category with small limits; duality supplies
the colimit statement when small colimits exist
\citep[Theorem 6.11]{Cisinski2003}. Consequently the bicomplete Quillen
case of Question 3.9 admits no counterexample of the requested kind.
For combinatorial models, both diagram model structures give a direct and
symmetric proof; \citet[Theorem 1.28 and Proposition 1.30]{Groth2013} presents
this route. The later existence theorem does not retroactively prove every
factorization assertion contemplated in the manuscript, and it does not
make the adjective ``very special'' necessary.

The final section shows why the size of the permitted diagrams matters.
Its bounded-complex example is effective precisely because bounded derived
categories need not contain the result of a homotopy limit indexed by a
finite category with arbitrarily long composable strings. For
$k=\mathbf F_2$ and $G=C_2$, the resolution of the trivial $k[G]$-module is
periodic and gives nonzero group homology and cohomology in every
nonnegative degree. Testing a hypothetical adjoint to the trivial-action
functor against shifts proves its nonexistence within the bounded category.
This supplies the required obstruction without identifying a bounded
category with a bicomplete modern model category. The distinction between
cofibration-category excision on cofibrant objects and the manuscript's
absolute $W$-cofibrations is equally necessary in reading the Brown lemma
invoked on manuscript page 35.

Some of the closing searches have elementary answers. A fully faithful
coreflective image in a bicomplete category is itself bicomplete: the
coreflector constructs limits, and the image is closed under colimits.
Thus the proposed adjoint-localization construction in Example 3.12 cannot
produce its hoped-for failure. The category of ordinals small relative to
the fixed universe has all small suprema and no terminal object, answering
Question 3.13 with the universe convention made explicit. Its lack of a
small conservative generating family is consistent with the convention of
SGA~4. A rationally ordered category cut at an irrational number supplies
a mutually lifting-closed pair which does not factor arrows across the cut;
it does not have the cocompleteness required of a special ambient category.

The broader localizer question in 3.14 admits a small counterexample.
Localize the Boolean lattice of subsets of $\{1,2,3\}$ at the single cover
$\{1\}\subset\{1,2\}$. The resulting seven-object skeleton has neither
a least upper bound for $\{2\},\{3\}$ nor a greatest lower bound for
$\{1,3\},\{2,3\}$. The localizer is strongly saturated, accessible, and
stable under filtered colimits, and the original lattice is complete and
cocomplete. The commentary computes the localization and proves each
assertion. This answers the broad question even with the additional
smallness hypotheses. The cofibration-category question in the same
paragraph has the opposite answer: the transfinite continuity axiom
supplies the coproduct and sequential-colimit hypotheses of modern
cofibration-category theory, which yields all homotopy left Kan
extensions \citep[Theorem 10.3.2]{RadulescuBanu2009}.

Question 3.15 also benefits from separating the closures introduced at the
start. The class of pure monomorphisms in directed multigraphs is accessible,
closed under filtered colimits, pushouts, transfinite composition, and
retracts. It nevertheless has no small set of Quillen generators. The
non-generation theorem is due to Cox, Feigert, Kamsma, Mazari-Armida, and
Rosick\'y, in their 2026 classification for presheaf categories
\citep[Theorem 5.1]{CoxEtAl2026Pure}; the
commentary verifies the translation into the hypotheses stated here and
supplies the small-subdiagram argument for accessibility. This answers the
question for a Q-saturated class. It does not identify that class with its
double lifting orthogonal, and the stronger question obtained by imposing
that additional condition is not silently substituted for the question
Grothendieck actually wrote.

\medskip
The commentary retains the edition's slips as source evidence, distinguishes
them from errors of transcription, and states repairs at the
point where they are needed. The claim register records the source's proof
status separately from later theorems or counterexamples. Mathematical
dependencies on Chapters XIII, XV, XVII, and XVIII are checked against
the relevant original passages or explicitly identified as recalled results.
\nocite{*}
\endgroup
